\section*{Chapter \ref{ch:s2:discretionary-macro} --- Discretionary Macro}

\begin{solution}{exo:s2:discretionary-macro:1}
$99\sqrt{0.5}/\sqrt{2\pi} = 27.9$ basis points.
\end{solution}

\begin{solution}{exo:s2:discretionary-macro:2}
The position with the 25 basis-point stop is twice the size of the one with the 50 basis-point stop: each loses one unit when stopped.
\end{solution}

\begin{solution}{exo:s2:discretionary-macro:3}
The spread pays its full width, 50 basis points: $50/18.2 - 1 = 1.75$ units of budget, net of the premium.
\end{solution}

\begin{solution}{exo:s2:discretionary-macro:4}
The chance that a path first touches a level falls quickly with the level's distance measured in the path's standard deviation: at six months the
rate's standard deviation is about 64 basis points, so a 10 basis-point stop is touched on most paths and a 100 basis-point one almost never.
\end{solution}

\begin{solution}{exo:s2:discretionary-macro:5}
It costs less (17.2 against 27.9 basis points), so the budget buys more of it and a large fall pays more per unit; but the rate must fall more than
25 basis points for it to pay at all, so it ends worthless more often.
\end{solution}

\begin{solution}{exo:s2:discretionary-macro:6}
Over six months the rate's own noise (about 64 basis points of standard deviation) is larger than the 50 basis-point fall expected, so on about 21\%
of paths the noise wins even when the drift is exactly as the view says.
\end{solution}

\begin{solution}{exo:s2:discretionary-macro:7}
With no dispersion the rate ends lower on 78.6\% of paths instead of 73.4\%, but most payoffs per unit of budget fall: the tight-stop futures from 1.68
to 1.45, the at-the-money option from 1.26 to 1.07, the out-of-the-money option from 1.69 to 1.32, the stop-50 futures from 0.98 to 0.91; only the
capped spread rises, from 0.71 to 0.79. Dispersion in the drift creates the large moves that convex expressions pay most for; precision helps
the expressions that give up large moves.
\end{solution}

\begin{solution}{exo:s2:discretionary-macro:8}
Being right about the direction is not the same as being right about the size, the timing and the path; and a year of outcomes is too few to tell skill
from noise. The trade journal should show whether the expression, the size or the stop turned a right view into a loss.
\end{solution}

\begin{solution}{pb:s2:discretionary-macro:1}
\begin{enumerate}
\item The instrument, strike, maturity and size that turn a view into a position.
\item Futures: the size; options: the size with a capped loss; stops: the path; spreads: a move up to a point.
\item A 50 basis-point fall in six months, drift dispersion 50, volatility 90 a year.
\item On 73.4\% of paths.
\item The premium one is willing to spend on options.
\item Options by their premium; stopped futures by their stop; unstopped futures by a two-standard-deviation loss.
\item 27.9, 17.2 and 18.2 basis points.
\item Each definition of risk scales positions differently.
\item A rule for cutting a position, set in advance and followed.
\item 55.6\%, 26.5\%, 5.6\%, 0.7\% and 0.1\% for 10, 25, 50, 75 and 100 basis points.
\item A framework for the value of stop-loss rules; some raise expected return while cutting volatility at longer frequencies.
\item A tighter stop means a larger position for the same budget.
\item \textbf{Named result.} Per unit of risk budget, futures earn 0.39, futures with a 25 or 50 basis-point stop 1.68 and 0.98, options 1.26 (at the
money) and 1.69 (out of the money) and the spread 0.71; a 25 basis-point stop cuts 26.5\% of correct views short, a 50 basis-point stop 5.6\%.
\item Size: futures or options; path: options rather than tight stops.
\item The view is right more often; convex expressions earn less.
\item The view, its horizon and confidence, the expression and its reason, and the outcome of each.
\item Separate the views' hit rate from the expressions' payoff, over many trades.
\item It is about expressing any view, not a strategy with a source of return.
\item Chapter~16 replaces the manager's views with rules.
\item The expression, the size and the stop.
\end{enumerate}
\end{solution}

\begin{solution}{iq:s2:discretionary-macro:1}
Receive fixed in short-dated swaps or buy the relevant futures (for example, SOFR futures for the meeting dates), or buy receiver options or calls on
the futures; choose by how sure one is of the timing and how much one can lose.
\end{solution}

\begin{solution}{iq:s2:discretionary-macro:2}
When the path is uncertain (the market may move against the view first), when the loss must be capped in advance, or when the view is about a large
move whose size is uncertain.
\end{solution}

\begin{solution}{iq:s2:discretionary-macro:3}
From the expected path's noise over the horizon: wide enough that noise rarely triggers them, sized so that a triggered stop costs a set share of the
book's loss limit; review them against the record of views stopped out.
\end{solution}

\begin{solution}{iq:s2:discretionary-macro:4}
Record views when formed and score them against outcomes over many trades, relative to a base rate and to a random expression; separate hit rate from
payoff per trade.
\end{solution}

\begin{solution}{iq:s2:discretionary-macro:5}
Each entry: time, view (direction, size, horizon, confidence), expression and size, stop and target, then exits with reason and P\&L; tagged so views
and expressions can be scored separately.
\end{solution}

\begin{solution}{iq:s2:discretionary-macro:6}
By the reflection principle, $P(\max_{t \le T} W_t \ge s) = 2\,(1 - N(s/(\sigma\sqrt T)))$. It falls fast once $s$ exceeds one standard deviation
$\sigma\sqrt T$ (about 64 basis points here), which is why the stop-out curve drops steeply between 10 and 50 basis points; the view's drift makes it
fall faster still.
\end{solution}
